% BHU PYQ -- Transcribed & Corrected
% Paper: CHEMJ33: Physical Chemistry-I
% Exam: B.Sc. (Hons) Semester III Examination, 2025-26 (NEP)
% Ref Code: CECONF/N/2025-26/057
% Category: Major
% Department: Chemistry

\documentclass[11pt,a4paper]{article}
\usepackage[a4paper,top=2.5cm,bottom=2.5cm,left=2.8cm,right=2.8cm,headheight=14pt]{geometry}
\usepackage{amsmath,amssymb}
\usepackage[T1]{fontenc}
\usepackage[utf8]{inputenc}
\usepackage{lmodern,microtype,fancyhdr,enumitem,hyperref,lastpage,parskip}

\hypersetup{
    colorlinks=true,
    urlcolor=black,
    linkcolor=black,
    pdftitle={CHEMJ33: Physical Chemistry-I -- BHU Sem III 2025-26 (NEP)}
}

\pagestyle{fancy}
\fancyhf{}
\renewcommand{\headrulewidth}{0.4pt}
\renewcommand{\footrulewidth}{0.4pt}
\fancyhead[L]{\small \textit{Banaras Hindu University}}
\fancyhead[R]{\small \textit{CHEMJ33: Physical Chemistry-I (2025-26)}}
\fancyfoot[C]{\small Page \thepage\ of \pageref{LastPage}}

\newcommand{\pts}[1]{\hfill \textit{[#1]}}

\begin{document}

\begin{center}
  {\large\bfseries BANARAS HINDU UNIVERSITY}\\[2pt]
  {\normalsize B.Sc. (Hons) Semester III Examination, 2025-26 (NEP)}\\[6pt]
  \rule{0.8\linewidth}{0.5pt}\\[6pt]
  {\large\bfseries CHEMISTRY}\\[4pt]
  {\large\bfseries Paper Code: CHEMJ33 : Physical Chemistry-I}\\[2pt]
  {\small Ref. No.: CECONF/N/2025-26/057}\\[4pt]
  {\normalsize Time: 2$\frac{1}{2}$ Hours\hfill Maximum Marks: 70}
\end{center}

\rule{\linewidth}{0.4pt}

\smallskip

\noindent \textit{Instructions: Answer five questions including Question 1 which is compulsory. The figures in the right-hand margin indicate marks.}

\bigskip

\noindent \textbf{Question 1.}
\begin{enumerate}[label=(\alph*)]
  \item Construct a plot for the conductometric titration of an aqueous mixture of HCl and $\text{CH}_3\text{COOH}$ solution with an aqueous sodium hydroxide solution and justify it. \pts{3}
  \item Write about the two types of electrodes and their electrode reactions. \pts{2}
  \item What is buffer solution? Write the conditions under which it can show maximum buffer-capacity. \pts{1+2}
  \item Discuss the eutectic point and the triple point with suitable examples. \pts{2+2}
  \item How can you determine the activation energy of a reaction? \pts{2}
\end{enumerate}

\bigskip

\noindent \textbf{Question 2.}
\begin{enumerate}[label=(\alph*)]
  \item Define molar conductance and equivalent conductance of an electrolyte. Write their units. \pts{3}
  \item Explain why proton ($\text{H}^+$) has a very high molar conductivity in water. \pts{3}
  \item What is the transference number of ion? Discuss the principle for the determination of the transference number using moving boundary method. \pts{1+3}
  \item Write the Kohlrausch's law of independent migration of ions. Discuss, with a suitable example, the use of it for the determination of the conductance of a weak electrolyte. \pts{2+2}
\end{enumerate}

\bigskip

\noindent \textbf{Question 3.}
\begin{enumerate}[label=(\alph*)]
  \item Write the relation between the ionic mobility and its limiting molar conductivity. Justify it. \pts{6}
  \item Discuss the applications of conductance measurements. \pts{3}
  \item Write the basic principle for the conductometric titrations. \pts{3}
  \item Predict the plot for the conductometric titration of $\text{NH}_4\text{OH}$ vs. HBr. \pts{2}
\end{enumerate}

\bigskip

\noindent \textbf{Question 4.}
\begin{enumerate}[label=(\alph*)]
  \item What is reference electrode? Discuss with a suitable example. \pts{2}
  \item Define the following terms: \pts{2+1}
  \begin{enumerate}[label=(\roman*)]
    \item Cell potential
    \item Standard electrode potential
  \end{enumerate}
  \item Write the relation between free energy change and cell potential of an electrochemical cell and justify it. From it, derive the Nernst equation for an electrochemical cell having a general cell reaction: $aA + bB \rightleftharpoons cC + dD$. \pts{3+3}
  \item Discuss the principle for the determination of the pH of a solution using a glass electrode. \pts{3}
\end{enumerate}

\bigskip

\noindent \textbf{Question 5.}
\begin{enumerate}[label=(\alph*)]
  \item Considering the following cell: $\text{Hg}(l) \mid \text{Hg}_2\text{Cl}_2(s) \mid \text{KCl (std. soln.)} \parallel \text{Fe}^{+2}(aq.), \text{Fe}^{+3}(aq.) \mid \text{Pt}(s)$ and assuming the redox reaction: $\text{Fe}^{+2}(aq.) + \text{Ce}^{+4}(aq.) = \text{Fe}^{+3}(aq.) + \text{Ce}^{+3}(aq.)$ involving the oxidation of $\text{Fe}^{+2}$ ions by $\text{Ce}^{+4}$ ions carried out in a potentiometric titration, write the expressions for the $E_{\text{cell}}$ before and after the equivalence point (Given: $E_{\text{calomel(satd.)}}^\circ = +0.242\text{ V}$, $E_{\text{Fe}^{+3}/\text{Fe}^{+2}}^\circ = +0.77\text{ V}$ and $E_{\text{Ce}^{+4}/\text{Ce}^{+3}}^\circ = +1.61\text{ V}$ at 298 K). \pts{3}
  \item Justify the nature of the expected plot to be obtained from the above titration. \pts{3}
  \item What is acid-base indicator? Derive a relation between pH a solution with the $pK_{\text{ind}}$ value of an acid indicator. What is the working pH range for the effective usefulness of an indicator? Justify your answer. \pts{2+2+2}
  \item Discuss the choice indicators for the titrations of (i) strong acid vs. strong base, (ii) weak acid vs. strong base, (iii) strong acid vs. weak base and (iv) weak acid vs. weak base. \pts{2}
\end{enumerate}

\bigskip

\noindent \textbf{Question 6.}
\begin{enumerate}[label=(\alph*)]
  \item Derive the Clapeyron-Clausius equation. \pts{3}
  \item Justify the phase rule. \pts{4}
  \item Use it to predict the maximum number of phases to be available for a single component system. Predict the position of these phases in the temperature vs. pressure phase diagram. \pts{3}
  \item Draw the phase diagram of water system at low pressure. Use the phase rule to calculate the no. of degrees of freedom (variance) at the following points (i) within the two curves, (ii) on the curve and (iii) at the junction points of all the curves. \pts{2+2}
\end{enumerate}

\bigskip

\noindent \textbf{Question 7.}
\begin{enumerate}[label=(\alph*)]
  \item Derive Nernst distribution law. \pts{2}
  \item Workout the distribution law of a solute undergoing dissociation in one solvent. \pts{2}
  \item Describe the procedure to determine the equilibrium constant of the reaction: $\text{KI} + \text{I}_2 \rightleftharpoons \text{KI}_3$ in aqueous medium using Nernst distribution law. \pts{4}
  \item Define the following terms: (i) Rate of reaction (ii) Rate laws of a reaction (iii) Order of a reaction. \pts{3}
  \item Derive the rate law for a zero-order reaction. Draw a kinetics plot for it. Give one example for such reaction. \pts{1+1+1}
\end{enumerate}

\bigskip

\noindent \textbf{Question 8.}
\begin{enumerate}[label=(\alph*)]
  \item Derive the integrated rate law for a first-order reaction considering the rate as the formation of the product. Write the characteristics of such reaction. \pts{2+1}
  \item Derive the integrated rate law for a second-order reaction considering the following equation: $A + B \rightarrow \text{Product}$. \pts{2}
  \item Workout the rate law for a pseudo first-order reaction. Justify that acid catalyzed inversion of cane sugar or hydrolysis of ester follows a pseudo first-order kinetics. \pts{2+1+2}
  \item Discuss the application of the steady state approximation concept with a suitable example. \pts{2}
  \item Write the three main factors affecting the rate of a reaction according to collision theory. \pts{2}
  \item Write the postulates of the transition state theory. \pts{2}
\end{enumerate}

\bigskip
\begin{center}
  \textbf{***}
\end{center}

\end{document}
